\section{チェビシェフ多項式}\label{節：チェビシェフ多項式}
\subsection{整数倍角とチェビシェフ多項式}

角度を2倍する式を見たので,\,一般の整数倍も考えよう.
多項式$T_m(x)$を
\[
 T_0(x)=1,\quad T_1(x)=x,
\]
\[
 T_{m+1}(x)=2xT_m(x)-T_{m-1}(x)
\]
で定める. これを\textbf{チェビシェフ多項式}という.
加法公式によって
\[
 T_m(\cos\theta)=\cos(m\theta)
\]
が帰納的に示される.
例えば
\[
 T_2(x)=2x^2-1,\qquad T_3(x)=4x^3-3x.
\]
特定の$x$を固定すると,$m$に関しては隣接3項間の線形漸化式である.

一方,\,固定した整数$d\ge2$について
\[
 a_{n+1}=T_d(a_n)
\]
を反復すると,$a_1=\cos\theta$のとき
\[
 a_n=\cos(d^{n-1}\theta)
\]
となる. こちらは$n$に関して非線形の更新である.
多項式の次数を進める添字$m$と,\,更新回数を表す$n$を区別しよう.

さらに
\[
 T_j(T_m(x))=T_{jm}(x)
\]
という合成の公式もある.
$x=\cos\theta$で確かめると両辺は$\cos(jm\theta)$になる.
両辺は多項式で,\,区間$[-1,1]$の無数の点で一致するため,\,すべての実数・複素数$x$で一致する.
零でない多項式の根が次数を超えないことを使っている.
従って一般の初項でも
\[
 a_n=T_{d^{n-1}}(a_1)
\]
と表せる.

\subsection{三角関数の外側にも続く式}

実数$a_1$が$[-1,1]$の外にあると,\,実角の余弦としては表せない.
しかし多項式の更新自体は定義されている.
任意の$x$に対し,\,非零の複素数$u$を
\[
 u^2-2xu+1=0
\]
の根として選べば,$x=\frac{u+u^{-1}}2$であり,
\[
 T_m(x)=\frac{u^m+u^{-m}}2
\]
となる. 初めの二つの値と更新を確かめればよい.
$|u|=1$なら余弦の表示に戻り,\,実数$x>1$なら実根を選んで指数的な増大として読める.

例えば$a_1=\frac54$,$a_{n+1}=2a_n^2-1$では$u=2$と選べるので,
\[
 a_n=\frac{2^{\,2^{n-1}}+2^{-2^{n-1}}}{2}.
\]
三角関数の公式は,\,二つの逆数の冪を足した式の一つの姿でもあった.
\chapterref{章：行列と連立漸化式}の等比的な成分への分解と,\,ここでも再会する.
